

\titre{Rattrapage METNU }
\theme{Méthode numérique}
\auteur{}
\organisation{AMSCC}

\contenu{






Soit \( f \) une application de \(\mathbb{R}\) dans \(\mathbb{R}\) définie par \( f(x) = e^{x^2} - 4x^2 \). On se propose de trouver les racines réelles de \( f \).

\begin{enumerate}
    \item Situer les quatre racines de \( f \), c'est-à-dire indiquer quatre intervalles disjoints qui contiennent chacun une et une seule racine.
    \item Montrer qu'il y a une racine \(\hat{x}\) comprise entre 0 et 1.
    \item Soit la méthode de point fixe : 
    \[
    x_0 \in ]0; 1[, \quad x_{k+1} = \varphi(x_k), \quad \text{où} \quad \varphi \text{ est l'application de } \mathbb{R} \text{ dans } \mathbb{R} \text{ définie par } \varphi(x) = \frac{1}{2}\sqrt{e^{x^2}}
    \]
    \begin{enumerate}
        \item Montrer qu'il existe une constante \( C \in ]0; 1[ \) telle que pour tout \( x \in [0; 1] \), \( |\varphi'(x)| < C \).
        \item Montrer que \( \varphi([0; 1]) \subset [0; 1] \).
        \item En déduire la convergence de la méthode proposée.
    \end{enumerate}
    \item On s'intéresse maintenant à la méthode de Newton
    \begin{enumerate}
        \item Écrire la suite \((x_k)_{k \in \mathbb{N}}\) des itérés de la méthode de Newton pour la recherche des zéros de \( f \).
        \item Donner un critère d'arrêt pour cette méthode.
    \end{enumerate}
\end{enumerate}
}